% In this section, we construct a model for the theory of \autoref{sec:type-theory} to show that it is consistent.

\subsection{Assemblies and partial equivalence relations}


In this realisability model, based on Hyland's effective topos, \cite{hyland_effective_1982}, types are interpreted as \emph{assemblies}: sets with associated computability data in the form of a relation between a set of \emph{realisers} and elements of the set.
The intuition given by Reus \cite{REUS1999128} is that realisers can be seen as `codes' in a low-level language or a general model of computation. A code that is related to an object in the set is said to realise it, and can be seen as a representation of that object. Maps between such assemblies must be computable in this general model of computation. 

The definition of an assembly is parameterised by a general model of computation, which is given by a partial combinatory algebra (pca) \cite{feferman}, providing the set of realisers. We will restrict ourselves to the most well-known pca, the first Kleene algebra \( \mathcal{K}_1 \), although the construction should work in general.
Here, the realisers are the natural numbers, which enumerate the partial recursive functions \(\{\varphi_n\}_{n \in \bN} \).  We define a partial operation \( \cdot : \bN \times \bN \rightharpoonup \bN \) on this set of realisers by \( m \cdot n := \varphi_m(n) \), and we write \( {m \cdot n \downarrow} \) when \( \varphi_m(n) \) is defined. 
\begin{definition}
    An \emph{assembly} is a pair $(X, \Vdash_X)$, where $X$ is a set and ${\Vdash_X} \subseteq \bN \times X$ is a relation such that for every \( x \in X \) there exists some \( n \in \bN \) such that \( n \Vdash_X x \). We will call \( \Vdash_X \) the \emph{realisability} relation of \( X \). A morphism between assemblies \((X, \Vdash_X)\) and \( (Y, \Vdash_Y) \) is a function \( f : X \to  Y \) such that there exists an \( e \in \bN \) (not part of the data) satisfying:
\[
    \forall n \in \bN.\,\forall x \in X.\,(n \Vdash_X x) \implies (e \cdot n \downarrow \mbox{ and } e \cdot n \Vdash_Y f(x)).
\] 
In this case, we say that \(f\) is \emph{tracked} by \(e\). We denote the category of assemblies by \( \Asm \).
\end{definition}
Note that there is a full and faithful inclusion from sets to assemblies defined by assigning the largest realisability relation; we denote this by \( \nabla X \coloneqq (X, \bN \times X) \). 

Since morphisms between assemblies are functions that are tracked by a realiser, they must be sufficiently computable.
This ensures that the dependent product can be small, even when quantifying over large types, which allows us to model the impredicative universal quantifier.
It also allows us to enforce a notion of parametricity in the model via the inclusion from sets to assemblies. The intuition is that \( \nabla X \) carries no meaningful computational data, making the objects of \( X \) indistinguishable when defining morphisms. Since we view the axioms that we added to the theory as a kind of parametricity, this construction allows us to enforce these axioms in the model. 

Since we will interpret types as assemblies, we need a sufficiently `small' subcategory to interpret the small types. As a step towards this goal, we first consider the following subcategory of assemblies.
\begin{definition}
An assembly $(X, \Vdash_x)$ is called a \emph{modest set} if every $n \in \bN$ codes at most one $a \in X$: 
\[
    \forall n \in \bN \. (n \Vdash_X x_1, x_2) \implies(x_1 = x_2).
\]
\end{definition}
We denote the category of modest sets by $\Mod$, which is a full subcategory of $\Asm$. However, the size of the carrier set of a modest set is not bounded in any way, so we cannot equip the `set' of all modest sets with the structure of an assembly.
Instead, we note that the elements of a modest set are not that important: we only need to know when two realisers encode the same element, as this lets us represent an object by the equivalence class of its realisers.
This insight gives us a category that is equivalent to the modest sets, namely that of partial equivalence relations on the set of realisers. 
\begin{definition}
    A \emph{partial equivalence relation (PER)}  is a symmetric and transitive relation $R \subseteq \mathbb{N} \times \mathbb{N}$. We define the following notation: 
\begin{align*}
	\rm{dom}(R) &:= \{n \in \mathbb{N} \mid n\mathbin Rn\}, &
	[n]_R &:= \{m \in \mathbb{N} \mid nRm\}, &
	\mathbb{N} /R &:= \{[n]_R \mid n \in \rm{dom}(R)\}.
\end{align*} 
A morphism between PERs R and S is a function \( f : \bN / R \to \bN /S \) such that \( f([n]_R) = [e \cdot n]_S \) for some \( e \in \bN \), in which case we say that \( f \) is tracked by \( e \).
We denote this category by \( \Per \).
\end{definition}
Note that a PER R is an equivalence relation on \( \rm{dom}(R) \).
\begin{lemma}\label{lemma:per-mod-equiv}
There is an equivalence between the categories $\Per$ and $\Mod$.
\end{lemma}
\begin{proof}[Proof Sketch]
    Given a PER \( R \subseteq \bN \times \bN \) the pair $\iota (R) \coloneqq  (\bN/R, \in)$ is a modest set.
    Conversely, given a modest set \( (X, \Vdash_X) \), we can define a relation on the set of realisers as follows: \begin{align*}
            n_1\sim_{X}n_2 \iff \exists x \in X \. (n_1 \Vdash_X x) \land (n_2 \Vdash_X x).
    \end{align*}
    This induces an equivalence between PERs and modest sets.
\end{proof}

In our model, we will use assemblies to interpret large types, while \autoref{lemma:per-mod-equiv} allows us to make use of the category of PERs to interpret the small types, since every PER induces a modest set with a countable carrier. 
To describe our model, we make use of the categories with families (CwF) framework, which closely follows the syntactic structure of dependent type theory.
% To show that a category \( \mathcal{C} \) can serve as a model for the structural rules of a dependent type theory, it suffices to show that \( \mathcal{C} \) can be given the structure of a CwF. 

We recall the basic parts of which a CwF \( \mathcal{C} \) consists:  
\begin{itemize}
    \item A category \( \C \) of \emph{semantic contexts} and context morphisms;
    \item For every semantic context \( \Gamma \in \C \), a collection \( \Ty(\Gamma) \) of \emph{semantic types};
    \item For every semantic context \( \Gamma \in \C \) and semantic type \( A \in \Ty(\Gamma) \), a collection \( \Tm(\Gamma, A) \) \emph{of semantic terms};
    \item For every context morphism $f : \Delta \to \Gamma$, an operation $-[f] : \Ty(\Gamma) \to \Ty(\Delta)$ for \emph{semantic substitution} on types, and for $A \in \Ty(\Gamma)$, an operation $-[f] : \Tm(\Gamma, A) \to \Tm(\Delta, A[f])$ for semantic substitution on terms;
    \item For every semantic context $\Gamma \in \C$ and $A \in \Ty(\Gamma)$, an extended context $\Gamma.A \in \C$. 
\end{itemize}

The category of assemblies can be given the structure of a CwF, with semantic contexts given by assemblies. A semantic type \( A \in \Ty(\Gamma) \) is given by a \( \Gamma \)-indexed family of assemblies \( (A_\gamma)_{\gamma \in \Gamma} \). Given a semantic type \( A \in \Ty(\Gamma) \), a semantic term \( a \in \Tm(\Gamma, A) \) is given by a \emph{dependent morphism} \( a : (\gamma \in \Gamma) \to A_\gamma \), that is, a dependent function tracked by an \(e\in\bb N\): 
\[
    \forall n \in \bN.\,\forall \gamma \in \Gamma \. (n\Vdash_\Gamma \gamma) \implies (e\cdot n\Vdash_{A_\gamma}a(\gamma)).
\]
We will also write this as a family \( (a_\gamma)_{\gamma \in \Gamma} \). The semantic substitution operation $-[f]$ on both types and terms is given by precomposing with $f$. Finally, context extension is given by dependent pairing (as in the ussual set-theoretic model).

We will omit the term `semantic', blurring the distinction between syntax and semantics. Further details on categories with families and the associated structures in assemblies can be found in, for example, \cite{Hofmann_1997, REUS1999128}. Assemblies support the interpretation of our basic types, dependent products and sums, and the identity type. In the following subsections, we highlight the relevant constructions.


\subsection{Universe type}
We give a brief description of how the universe type is interpreted in the model, as it will provide intuition on why we interpret the type of sizes in the same way.
The universe type \( U \) in a context \( \Gamma \) is interpreted as a semantic type \( U \in \Ty(\Gamma) \), meaning a \( \Gamma \)-indexed family of assemblies. Since the universe type is a non-dependent type, its interpretation should be a constant family. As we will want to interpret the small types using PERs, we will interpret \( U \) as the constant family of an assembly where the carrier set is the set of all PERs, which we also denote by \( \Per \). To construct an assembly out of \( \Per \), we use the inclusion from sets to assemblies, giving us \( \nabla \Per \). Recall that the intuition is that \( \nabla \Per \) carries no meaningful computational data, making the PERs computationally indistinguishable when defining morphisms.
The decoder \( \El \) is interpreted by the inclusion of PERs into assemblies.
% The decoder \( \El \) is interpreted as a semantic type in the extended context \( \Gamma.U \), meaning \( \El \in \Ty(\Gamma.U) \). Thus, we define it as a \( \Gamma.U \)-indexed family of assemblies where \( \El_{(\gamma, R)} = \iota (R) \), noting that \( R \in \nabla \Per \) is a PER and \( \iota  \) is the inclusion of PERs into assemblies. Given a small type \( A \in \Tm(\Gamma, U) \), we obtain a type \( \El[\overline{A}] \in \Ty(\Gamma) \), which is defined by \( \El[\overline{A}]_\gamma = \iota(A_\gamma) \).

\subsection{Type of sizes}
In previous work \cite{parametric-quantifier}, the type of sizes is interpreted as the set of natural numbers. By contrast, the impredicativity of our realisability model allows us to interpret the type $\Size$ as a large type, while still having universal quantification over $\Size$ with a small type live in $\U$. Given a context $\Gamma$, we define the semantic type $\Size \in \Ty(\Gamma)$ as $\Size_\gamma = \nabla \omega_1$, where $\omega_1$ is the first uncountable ordinal: the set of all countable ordinals.
We construct an assembly from the ordinal \( \omega_1 \) by giving it the full realisability relation, just as we did for the interpretation of the universe type. This ensures that morphisms from \( \nabla \omega_1 \) must treat all sizes uniformly.

The reason for interpreting the type of sizes as $\omega_1$ is that it allows us to justify the following equivalence:
\begin{align*}
	\Pi{x : A} \. \exists i . B(x, i)&\simeq\exists i\. \Pi {x : A} \. \exists j < i . B(x, j). \tag{\(\exists\) parametric over \(\Pi\)}
\end{align*}
Intuitively, the left-to-right direction of this equivalence states the following: given a dependent function that sends a term \( a : A \) to a pair $\langle s, b \rangle : \exists i . B(a, i)$, there is a limit size $\alpha$ such that $s < \alpha$ for each $\brac{s, b}$ in the range of the dependent function. This requires limits of sequences of sizes to be a size as well. Note that, since $A$ here is a small type, it is interpreted as a partial equivalence relation over $\bN$, hence it has countably many equivalence classes. Since countable limits of countable ordinals are themselves countable, $\omega_1$ suffices for our purpose.

The ordering on sizes is interpreted as the ordering on \( \omega_1 \), interpreting \( s_1 \le s_2 \) as the terminal singleton assembly when true, and the initial empty assembly otherwise. This interpretation validates the encoding of the ordering on sizes as a proposition in the theory.

\subsection{Bounded quantification and the fixpoint operator}
Recall that the notation $i : \Size \vdash \forall j < i \. A(j)$ for bounded universal quantification is shorthand for $i : \Size \vdash \forall j \. (j < i \to A(j))$. A term of this type is interpreted by a family of dependent morphisms, where the domain is bounded by the size in context, which confirms the syntactic intuition for this type.

%Thus, a term $g$ of type $\forall j < i \. A(j)$ is interpreted as a $\Gamma.\Size$-indexed collection of dependent morphisms, where
%\[
%g_{(\gamma, \alpha)} : (\beta \in \Size_\gamma) \to ((\beta < \alpha)_\gamma \to A_{(\gamma, \beta)})
%\]
%
%If $(\beta < \alpha)_\gamma $ holds in the interpretation, then $(\beta < \alpha)_\gamma$ is the terminal singleton assembly. and thus the set of morphisms of type $(\beta < \alpha)_\gamma \to A_{(\gamma, \beta)}$ is isomorphic to $A_{(\gamma, \beta)}$, allowing us to treat $g_{(\gamma, \alpha)}(\beta)$ as an element of $A_{(\gamma, \beta)}$. 
%If $(\beta < \alpha)_\gamma$ does not hold in the interpretation, then $(\beta < \alpha)_\gamma$ is the initial (empty) assembly and there thus exists a unique morphism of type $(\beta < \alpha)_\gamma \to A_{(\gamma, \beta)}$, namely the empty morphism, again given by $g_{(\gamma, \alpha)}(\beta)$. This confirms the syntactic intuition that we may treat $g_{(\gamma, \alpha)}$ as a dependent morphism where $\dom(g_{(\gamma, \alpha)}) = \Size_\gamma \upharpoonright \alpha$. We will thus write the type of $g_{(\gamma, \alpha)}$ as:
%\[
%g_{(\gamma, \alpha)} : (\beta \in \Size_{<\alpha, \gamma}) \to A_{(\gamma, \beta)}
%\]

Given a term \( f : \forall i. (\forall j < i. A(j)) \to A(i) \), the fixpoint operator \( \fix \, f : \forall i . A(i) \) is interpreted as a morphism, where the underlying dependent function is defined by well-founded induction on \( \omega_1 \). The definition given by well-founded induction corresponds exactly to the \( \beta \)-rule in the theory. In order for \( \fix \, f \) to be a morphism, it must be tracked by a realiser. The following fact, which states that a fixpoint operator exists in every partial combinatory algebra, allows us to find an appropriate realiser for the fixpoint operator.
\begin{proposition}[Fixpoint operator in a pca \cite{oosten}]
    For every pca $(\cA, \cdot)$, there exists a term $\fix\! \in \cA$ such that the following hold for any \( f, a \in \cA \) we have $\fix\,f \downarrow$ and $\fix\,f \, a \simeq f \, (\fix\,f) \, a$.
\end{proposition}


\subsection{Universal and existential type}
The assembly model supports the interpretation of general impredicative products. Although we do not add impredicative products to the theory in full generality, this does allow us to interpret the (impredicative) universal quantifier $\forall i . A(i)$ where $A : U$ in the model simply as a dependent product. The interpretation of sizes as $\nabla \omega_1$ then suffices to validate the equivalences for the universal quantifier.

To interpret the existential type \( \exists i.A(i) \), we construct a quotient of the corresponding \( \Sigma \)-type that lives inside the universe. Thus, we should canonically construct a family of PERs from the semantic $\Sigma$-type, \( \Sigma(\Size, A) \), which is a family of assemblies defined in \autoref{appendix:model}. The following lemma gives us a quotient construction resulting in a modest set.
\begin{lemma}
    The inclusion functor \( I : \Mod \to \Asm \) has a left adjoint \( \M : \Asm \to \Mod \).
\end{lemma}
\begin{proof}[Proof Sketch]
    For an assembly $(A, \Vdash_A)$, we define a modest set $\M(A) = ( A {/ \!\sim_\M}, \Vdash_{\M(A)})$ where $\sim_\M$ is the least equivalence relation satisfying the following: $\exists n \in \bN \. (n \Vdash_A a_1, a_2) \implies (a_1 \sim_\M a_2)$.
The relation is defined by
\(
n \Vdash_{\M(A)} [ a ]_{\sim_\M} \Iff \exists a' \in [ a ]_{\sim_\M} (n \Vdash_A a').\)
\end{proof}
However, the universe consists of PERs rather than modest sets. Following \autoref{lemma:per-mod-equiv}, every modest set is isomorphic to a PER. Thus, we may treat \( \M(A) \) as the PER that is isomorphic to it. Given a type \( A \in \Ty(\Gamma) \), we can define its \(U\)-\emph{truncation} \( \Tr{A} \in \Tm(\Gamma, U) \) as the family of PERs \( \Tr{A}_\gamma = \M(A_\gamma) \). 
Given a small type \( A \in \Tm(\Gamma, U) \), this allows us to interpret the small type \( \exists i. A(i) \in \Tm(\Gamma, U)\) as the \(U\)-truncation of the \( \Sigma \)-type: 
\[
    \exists i. A(i) \coloneqq \Tr{\Sigma(\Size, \El A)}.
\] 

This also allows us to interpret the pairing operation and elimination of existential types. Moreover, it validates the equivalences we have axiomatised in the theory related to the existential quantifier in the model. Further details of the model are given in \autoref{appendix:model}.
